\begin{abstract}
\vspace{0.25in}
Sparse, delayed, or deceptive rewards remain a persistent challenge in deep reinforcement learning because a policy may receive little task-relevant feedback before useful behavior is discovered.
Intrinsic motivation addresses this limitation by augmenting the environment reward with auxiliary signals that encourage exploration, but common novelty- and prediction-error-based objectives can reward transitions that are surprising without being learnable or useful.
This study investigates an intrinsic-reward method for reducing such unproductive curiosity while preserving the exploratory benefits of reward augmentation.
The proposed method, \emph{Gated Learning-Progress Exploration} (GLPE), combines two complementary transition-level signals: feature-space impact, measured by the change between consecutive latent observations, and region-local learning progress, measured by the decrease in forward-dynamics prediction error within an online partition of the learned latent space.
The intrinsic reward is computed from normalized impact and learning-progress components and is optionally filtered by a region-specific gate that suppresses rewards in areas with persistently high prediction error but low learning progress.
An ungated variant is also evaluated to isolate the effect of the base intrinsic score from the gate.

The method is implemented with a Proximal Policy Optimization backbone and evaluated on five Gymnasium control benchmarks: \textsc{MountainCar-v0}, \textsc{BipedalWalker-v3}, \textsc{Ant-v5}, \textsc{HalfCheetah-v5}, and \textsc{Humanoid-v5}.
The experimental design compares GLPE and GLPE without gating against PPO, ICM, RND, RIDE, and RIAC under shared policy architectures, training budgets, and evaluation procedures.
Across the benchmark suite, GLPE without gating is a consistent default that remains competitive with the strongest baseline methods under both sample-efficiency and wall-clock views.
The gated variant is most useful in sparse-reward settings such as \textsc{MountainCar-v0}, where it achieves complete solved-threshold reliability while producing the strongest final return and step-AUC, but it can be conservative in dense-reward or high-variance tasks.
Ablation results show that impact and learning progress contribute differently across environments, supporting their combined use as a robust exploration signal.
The study contributes a formal formulation of GLPE, an empirical comparison against established intrinsic-motivation baselines, and an analysis of the tradeoff between targeted exploration control and computational overhead.
\end{abstract}